\subsection{同伦等价的空间}

定义 2 —— 设 X 与 Y 是拓扑空间。同伦类 \(\varphi \in [\mathrm{X};\mathrm{Y}]\) 称为可逆的，如果存在同伦类 \(\psi \in [\mathrm{Y};\mathrm{X}]\)，使得 \(\psi \circ \varphi = [\mathrm{Id}_{\mathrm{X}}]\) 且 \(\varphi \circ \psi = [\mathrm{Id}_{\mathrm{Y}}]\)。若一个连续映射的同伦类可逆，则称该连续映射为同伦等价。

设 \(\varphi\in[X;Y]\) 是一个可逆同伦类。具有定义 2 中诸性质的同伦类 \(\psi\in[Y;X]\) 是唯一的。事实上，设 \(\psi,\psi'\) 是具有这些性质的类；我们有

\[
\psi = \psi \circ [ \mathrm{Id} _ {\mathrm{Y}} ] = \psi \circ \varphi \circ \psi^ {\prime} = [ \mathrm{Id} _ {\mathrm{X}} ] \circ \psi^ {\prime} = \psi^ {\prime}.
\]

这个唯一的类称为同伦类 \(\varphi\) 的逆，记作 \(\varphi^{-1}\)。

设 Z 是拓扑空间。我们有 \(\chi\circ\varphi\circ\varphi^{-1}=\chi\)（对所有 \(\chi\in[Y;Z]\)），且 \(\theta\circ\varphi^{-1}\circ\varphi=\theta\)（对所有 \(\theta\in[X;Z].\)）。由此推出，映射 \(\chi\mapsto\chi\circ\varphi\)（从 {[}Y;Z{]} 到 {[}X;Z{]}）与映射 \(\theta\mapsto\theta\circ\varphi^{-1}\)（从 {[}X;Z{]} 到 {[}Y;Z{]}）互为逆双射。同样，映射 \(\chi\mapsto\varphi\circ\chi\)（从 {[}Z;X{]} 到 {[}Z;Y{]}）是双射，其逆是映射 \(\theta\mapsto\varphi^{-1}\circ\theta\)（从 {[}Z;Y{]} 到 {[}Z;X{]}）。

设 X、Y、Z 是拓扑空间，\(\varphi \in [X;Y]\) 与 \(\psi \in [Y;Z]\) 是可逆同伦类。于是类 \(\psi \circ \varphi\) 可逆，其逆为 \(\varphi^{-1} \circ \psi^{-1}\)。事实上，我们有

\[
\begin{array}{r l} & (\psi \circ \varphi) \circ (\varphi^ {- 1} \circ \psi^ {- 1}) = \psi \circ (\varphi \circ \varphi^ {- 1}) \circ \psi^ {- 1} \\ & \qquad = \psi \circ [ \mathrm{Id} _ {\mathrm{Y}} ] \circ \psi^ {- 1} \\ & \qquad = \psi \circ \psi^ {- 1} = [ \mathrm{Id} _ {\mathrm{Z}} ] \end{array}
\]

类似地，\((\varphi^{-1} \circ \psi^{-1}) \circ (\psi \circ \varphi) = [\mathrm{Id}_{\mathrm{X}}]\)。

设 X 与 Y 是拓扑空间，\(f: X \to Y\) 是一个同伦等价。凡其类为 f 的类之逆的连续映射 \(g: Y \to X\)，均称为 f 的同伦逆（或逆）。这样的映射 g 是一个同伦等价。

同胚 \(f\) 是同伦等价，且 \([f]^{-1} = [f^{-1}]\)。两个同伦等价的复合是同伦等价。关系“X 与 Y 是拓扑空间且存在从 X 到 Y 的同伦等价”是一个等价关系。

定义 3 —— 若存在从 X 到 Y 的同伦等价，则称拓扑空间 X 与 Y 同伦等价。

例 —— 1) 空拓扑空间只与它自身同伦等价。

\begin{enumerate}
\def\labelenumi{\arabic{enumi})}
\setcounter{enumi}{1}
\item
  设 X 是非空拓扑空间。X 与一个化为一点的空间同伦等价的充要条件是：存在与 X 的恒同映射同伦的常值映射 \(p: X \rightarrow X\)。（此时对每个常值映射 \(q \colon X \to X\) 亦然，因为 \([p] = [p] \circ [q] = [\mathrm{Id}_{\mathrm{X}}] \circ [q] = [q]\)。）事实上，设 \(P\) 是化为一点的空间，\(f \colon P \to X\) 是一个映射，\(g \colon X \to P\) 是从 X 到 P 的唯一映射。我们有 \(g \circ f = \mathrm{Id}_{\mathrm{P}}\)；f 是同伦等价的充要条件是 \(f \circ g\) 与 \(\mathrm{Id}_{\mathrm{X}}\) 同伦。而 \(f \circ g\) 是常值的，其像为 \(f(P)\)。
\item
  若存在在 x 处带基点的同伦，把 \(Id_{X}\) 连接到从 X 到 X、以 \(\{x\}\) 为像的常值映射，则称带基点拓扑空间 (X,x) 是可缩的，或称拓扑空间 X 在 x 处可缩。这样的空间与一点同伦等价。然而，存在与一点同伦等价但在其任何点处都不可缩的空间，也存在在某一点处可缩、但并非在每一点处都可缩的空间（III, p. 321, 习题 1）。
\item
  设 E 是 n 维欧氏空间（或更一般地，R 上的拓扑向量空间），X 是 E 的一个子集。若对每个 \(y \in X\) 及每个 \(t \in I\)，点 \(tx + (1 - t)y\) 都属于 X，则称集合 X 在其一点 x 处是星形的。E 的凸子集（I, p. 122）在其每一点处都是星形的。
\end{enumerate}

在其一点 x 处星形的 E 的拓扑子空间 X 在该点处可缩。事实上，映射 \(\sigma\colon X \times I \to X\)（由 \(\sigma(y,t) = tx + (1 - t)y\) 定义）是在 x 处带基点的同伦，它把 \(Id_{X}\) 连接到以 \(\{x\}\) 为像的常值映射。

特别地，R 的每个区间、欧氏空间的每个凸子集，或更一般地，R 上拓扑向量空间的每个凸子集，都在其每一点处可缩。

\begin{enumerate}
\def\labelenumi{\arabic{enumi})}
\setcounter{enumi}{4}
\tightlist
\item
  我们将在后面（TA, V）证明，欧氏球面 \(S_n\)（\(n\geqslant1\)）不与一点同伦等价。可分无穷维希尔伯特空间的单位球面在其每一点处都可缩（EVT, V, p. 71, 习题 13）。
\end{enumerate}

